ANIMA Phase I investigated emergent organization and the architectural
limits of persistent learning in a minimally specified spiking nervous
system. Under frozen protocols, byte-identical identity gates, and
preserved negative results, the record demonstrates that:

\begin{itemize}
\item sensory-specific synaptic organization emerges from local
plasticity without supervision;
\item transient spike-count representations exist but separate poorly
across patterns;
\item every persistent substrate tested --- the slow intrinsic state
and the synaptic weight structure alike --- fails to retain episode
identity through repeated or alternated exposure, in the testable case
already during the stimulus;
\item bounded consolidation produces stable protected structure and
preserves alternating coexistence, but blocked-order allocation
remains recency-dominated;
\item the allocation sequence localized the residual failure to the
working-substrate churn that precedes a never-seen pattern's first
exposure, and from there to the recruitment of postsynaptic activity
for genuinely novel late-arriving patterns.
\end{itemize}

Phase I closes with that capability unresolved. The final experiment
tested the smallest hypothesized local amendment --- a fixed-magnitude,
input-proportional recruitment gain --- and was closed by its
pre-registered rules at two registrations, having established the
operating-point failure mode and the structural bound (a gain cannot
amplify a substrate destroyed before exposure), without either
destroying or confirming the weak-pattern recruitment hypothesis for
the class of mechanisms. What remains is not a tuning problem within
the tested architecture but an architectural question: how a substrate
whose persistent store is a shared additive budget can form multiple
independently retrievable persistent representations under sequential
experience.